%----------------------------------------------------------------------------------------
%	Ejercicios Capítulo 10
%----------------------------------------------------------------------------------------

\chapterimage{A_Capitulo_Ejercicios/img/portada/ima2} % Chapter heading image

\chapter*{Capítulo 10}
\addcontentsline{toc}{chapter}{\textcolor{ocre}{Capítulo 10}}

\begin{itemize}

 \item 1) Calcule el CAUE de una máquina que tiene un costo inicial de COP 8.000 y un valor de salvamento de COP 500 después de 8 años. Los costos anuales de operación (CAO) para la máquina se estiman en COP 900 y la tasa de interés es 20\% anual.\\
       \textbf{Respuesta:} COP 2.955/año (exacto: COP 2.954,57/año).
       \medskip

 \item 2) Los siguientes costos son los estimados para dos máquinas peladoras de tomates en una fábrica de conservas:

       \begin{center}
        \renewcommand{\arraystretch}{1.2}
        \begin{tabular}{|l|c|c|}
         \hline
         \rowcolor[HTML]{FFB183}
         \textbf{Factor} & \textbf{Máquina A} & \textbf{Máquina B} \\ \hline
         Costo inicial                & COP 26.000         & COP 36.000         \\ \hline
         Costo anual de mantenimiento & COP 800            & COP 300            \\ \hline
         Costo anual de mano de obra  & COP 11.000         & COP 7.000          \\ \hline
         Ingresos adicionales         & COP 2.000          & COP 3.000          \\ \hline
         Vida útil (años)             & 6                  & 10                 \\ \hline
        \end{tabular}
       \end{center}

       Si la tasa de retorno mínima requerida es 15\% anual, ¿qué máquina debe seleccionarse?\\
       \textbf{Respuesta:} Se selecciona la máquina B puesto que $CAUE_B$ (COP 16.924,53/año) $<$ $CAUE_A$ (COP 18.442,16/año).
       \medskip

 \item 3) Si un inversionista deposita COP 10.000 hoy, a una tasa de interés de 7\% anual, ¿cuántos años deberá acumular el dinero antes de que pueda retirar COP 1.400 anuales indefinidamente?\\
       \textbf{Respuesta:} 10,24 años.
       \medskip

 \item 4) Utilizando los valores del ejercicio 1, calcule el CAUE usando el método de la recuperación del capital más costo anual de mantenimiento.\\
       \textbf{Respuesta:} COP 2.955/año (exacto: COP 2.954,57/año).
       \medskip

 \item 5) Utilizando los valores del ejercicio 1, calcule el CAUE usando el método de la recuperación del capital únicamente.\\
       \textbf{Respuesta:} COP 2.955/año (exacto: COP 2.954,57/año).
       \medskip

\end{itemize}
